% =====================================================================
\section{Органы управления: энкодер, клавиатура, сенсорные кнопки}
% =====================================================================
\lessonintro
  {подключать три распространённых способа ввода и выбирать подходящий для своего устройства}
  {урок 3 (кнопка, подтяжка, дребезг, прерывания), урок 5 (потенциометр)}
  {поворот ручки меняет счётчик, клавиатура печатает нажатую клавишу, касание провода зажигает светодиод}
  {модуль энкодера KY-040, матричная клавиатура 4×4, отрезок провода или фольги}

\subsection{Сначала: как человек разговаривает с устройством}

Одной кнопки хватает, чтобы включить и выключить. Но чтобы выбрать пункт меню, набрать код или
плавно настроить громкость, нужны другие органы управления. В этом уроке их три. Энкодер --- ручка,
которая крутится бесконечно и сообщает, в какую сторону её повернули. Матричная клавиатура ---
шестнадцать кнопок, которым хватает восьми выводов. Сенсорная кнопка --- контакт без механики, который
чувствует прикосновение пальца.

\begin{center}\small
\begin{tabularx}{\textwidth}{@{}l X X@{}}
\toprule
Орган управления & Удобен для & Особенность \\
\midrule
Кнопка & включить, подтвердить & один вывод на кнопку \\
Потенциометр & плавная настройка с упорами & помнит положение; нужен АЦП \\
Энкодер & меню, громкость, счёт & крутится бесконечно; положение считает программа \\
Клавиатура 4×4 & ввод чисел и кодов & 16 кнопок на 8 выводах \\
Сенсорная кнопка & панель под стеклом или пластиком & нет движущихся частей; нужна калибровка \\
\bottomrule
\end{tabularx}
\end{center}

\subsection{Энкодер}

Внутри энкодера два контакта, A и B, которые при вращении замыкаются со сдвигом. По тому, какой из
них сработал первым, определяют направление: в момент спада сигнала A смотрят на уровень B.

\begin{figure}[H]
\centering
\begin{tikzpicture}[x=0.76cm,y=0.8cm,font=\sffamily\scriptsize]
  \node[left,text=accent] at (0,3.4) {A (CLK)};
  \node[left,text=okgreen] at (0,1.4) {B (DT)};
  \draw[very thick,accent] (0,4) -- (2,4) -- (2,3) -- (4,3) -- (4,4) -- (6,4) -- (6,3) -- (8,3) -- (8,4) -- (9,4);
  \draw[very thick,okgreen] (0,2) -- (3,2) -- (3,1) -- (5,1) -- (5,2) -- (7,2) -- (7,1) -- (9,1);
  \foreach \x in {2,6}{
    \draw[dashed,esp] (\x,0.6) -- (\x,4.5);
    \fill[esp] (\x,2) circle (2.2pt);
  }
  \node[above,text=esp,align=center] at (4,4.5) {спад A, а B ещё HIGH};
  \node[below] at (4.5,0.5) {поворот по часовой стрелке: $+1$};
  \begin{scope}[shift={(10.5,0)}]
    \draw[very thick,accent] (0,4) -- (3,4) -- (3,3) -- (5,3) -- (5,4) -- (7,4) -- (7,3) -- (9,3);
    \draw[very thick,okgreen] (0,2) -- (2,2) -- (2,1) -- (4,1) -- (4,2) -- (6,2) -- (6,1) -- (8,1) -- (8,2) -- (9,2);
    \foreach \x in {3,7}{
      \draw[dashed,esp] (\x,0.6) -- (\x,4.5);
      \fill[esp] (\x,1) circle (2.2pt);
    }
    \node[above,text=esp,align=center] at (5,4.5) {спад A, а B уже LOW};
    \node[below] at (4.5,0.5) {поворот против часовой: $-1$};
  \end{scope}
\end{tikzpicture}
\caption{Сигналы энкодера. Направление вращения определяет уровень B в момент спада A}
\end{figure}

Модуль KY-040 подключают так: «+» к 3V3, GND к земле, CLK к \pin{6}, DT к \pin{7}, SW (кнопка на оси)
к \pin{15}. Спад сигнала A ловит прерывание, поэтому ни один щелчок не теряется, даже если основная
программа занята.

\codecap{Энкодер: счёт поворотов и сброс нажатием на ось}
\codeinput{l23_encoder}

\subsection{Матричная клавиатура}

Шестнадцать кнопок клавиатуры расположены на пересечениях четырёх строк и четырёх столбцов. Кнопка
соединяет «свою» строку со «своим» столбцом.

\begin{figure}[H]
\centering
\begin{tikzpicture}[x=1.5cm,y=0.95cm,font=\sffamily\scriptsize]
  \foreach \r/\p in {0/15,1/16,2/17,3/18}{
    \draw[thick,accent] (-0.6,-\r) -- (3.6,-\r);
    \node[left,text=accent] at (-0.6,-\r) {строка \r\ (GPIO\p)};
  }
  \foreach \c/\p in {0/6,1/7,2/21,3/47}{
    \draw[thick,okgreen] (\c+0.3,0.75) -- (\c+0.3,-3.6);
    \node[above,text=okgreen,align=center] at (\c+0.3,0.75) {столбец \c\\ GPIO\p};
  }
  \foreach \r/\row in {0/{1,2,3,A},1/{4,5,6,B},2/{7,8,9,C},3/{*,0,\#,D}}{
    \foreach \k [count=\c from 0] in \row{
      \node[draw,circle,fill=white,inner sep=1.2pt,minimum size=5.5mm,font=\ttfamily\small] at (\c,-\r-0.32) {\k};
    }
  }
  \node[anchor=west,align=left] at (3.9,-0.6) {1. Строку 1 опускаем в LOW,\\ \phantom{1. }остальные отключаем.};
  \node[anchor=west,align=left] at (3.9,-1.7) {2. Читаем четыре столбца\\ \phantom{2. }(на них подтяжка к HIGH).};
  \node[anchor=west,align=left] at (3.9,-2.8) {3. Столбец 2 стал LOW ---\\ \phantom{3. }нажата клавиша «6».};
\end{tikzpicture}
\caption{Матрица 4×4: кнопка стоит на пересечении строки и столбца. Строки опрашиваются по очереди}
\end{figure}

Программа по очереди делает активной (LOW) одну строку и читает столбцы. Неактивные строки
переводятся в режим входа, то есть как бы отключаются: иначе при двух нажатых кнопках два выхода
с разными уровнями окажутся соединены накоротко.

\codecap{Опрос матричной клавиатуры 4×4}
\codeinput{l23_keypad}

\subsection{Сенсорные входы}

Выводы \pin{1}--\pin{14} умеют измерять электрическую ёмкость. Палец увеличивает ёмкость контакта, и число,
которое возвращает \code{touchRead()}, растёт. Достаточно вставить в \pin{14} отрезок провода или
приклеить к нему кусочек фольги.

\codecap{Сенсорная кнопка с калибровкой при запуске}
\codeinput{l23_touch}

Порог нельзя задать раз и навсегда: уровень зависит от длины провода, влажности и даже от того,
питается ли плата от ноутбука или от батареи. Поэтому программа измеряет «спокойный» уровень при
запуске и ставит порог чуть выше.

\begin{caution}
На первом ESP32 значение \code{touchRead()} при касании \emph{уменьшается}, а на ESP32-S3 ---
\emph{увеличивается}. Примеры из интернета для старого чипа нужно переписывать с обратным условием.
\end{caution}

\subsection{Самостоятельный опыт: меню на энкодере}

Соедините энкодер с OLED-дисплеем из урока 6. На экране --- три строки меню («Яркость», «Звук», «Выход»),
поворот ручки перемещает указатель, нажатие на ось печатает выбранный пункт в Serial Monitor. Указатель
не должен уходить за первую и последнюю строку.

\begin{trouble}
\begin{itemize}[leftmargin=*]
  \item \textbf{Энкодер считает через раз или «скачет»} --- дребезг контактов: увеличьте паузу в обработчике
        или поставьте конденсаторы \SI{100}{nF} между CLK, DT и GND.
  \item \textbf{Клавиатура выдаёт не те символы} --- шлейф подключён в обратном порядке; поменяйте порядок
        в массивах \code{ROWS} и \code{COLS}.
  \item \textbf{Сенсор срабатывает сам} --- провод слишком длинный или лежит рядом с источником помех;
        укоротите его и увеличьте запас порога.
\end{itemize}
\end{trouble}

\begin{summary}
\begin{itemize}[leftmargin=*]
  \item Энкодер сообщает направление: в момент спада A читаем уровень B.
  \item Матричная клавиатура опрашивается по строкам; неактивные строки переводим в режим входа.
  \item Сенсорный вход измеряет ёмкость; порог срабатывания калибруют при запуске.
\end{itemize}
\end{summary}

\begin{selfcheck}
\begin{enumerate}
  \item Как по двум сигналам энкодера определить направление вращения?
  \item Почему 16 кнопкам матричной клавиатуры хватает восьми выводов?
  \item Зачем сенсорной кнопке калибровка при запуске?
\end{enumerate}
\end{selfcheck}

\begin{tasks}
\begin{enumerate}
\item \lv{1} Запустите пример с энкодером и выведите в Serial Monitor направление вращения.
\item \lv{1} Определите нажатую клавишу матричной клавиатуры и выведите её в Serial Monitor.
\item \lv{2} Управляйте яркостью светодиода поворотом энкодера: один щелчок --- 5\%.
\item \lv{2} Выведите на экран вводимые цифры с клавиатуры, клавиша «\#» подтверждает ввод.
\item \lv{3} Сделайте кодовый замок: четыре цифры и «\#»; верный код зажигает зелёный светодиод, неверный --- красный.
\item \lv{4} Сделайте сенсорный «слайдер» из трёх контактов и определяйте направление движения пальца.
\end{enumerate}
\end{tasks}
